\section{Introduction}
\label{sec:introduction}

Low Earth orbit is becoming a setting for increasingly physical spacecraft
behaviors: robotic inspection, servicing, deployable structures, on-orbit
assembly, capture, docking, and contact with targets whose mass properties may
be uncertain. These tasks are not purely orbital, purely attitude, or purely
robotic. They require a simulator that can represent a spacecraft as an
articulated multibody system while also applying the environmental and
actuator loads that dominate orbital flight.

Astrodynamics simulators typically provide accurate orbit propagation,
environment models, and flight-dynamics workflows, but they are not designed
around general contact-rich articulated dynamics. Robotics simulators provide
fast rigid-body dynamics, constraints, actuation, and contact, but usually
assume a terrestrial inertial or gravity-fixed world. \mjorbit{} is motivated
by the gap between these two tool families. The package treats \mj{} as the
inner multibody dynamics engine and adds an orbital layer that keeps the
MuJoCo world frame aligned with a local LVLH frame attached to a propagated
chief orbit.

The current implementation exposes a compact public API:
\texttt{MjoModel} stores the compiled MuJoCo model and static orbital-coupling
metadata, \texttt{MjoData} stores one simulation rollout, and
\texttt{mjo\_forward}/\texttt{mjo\_step} mirror MuJoCo's synchronization and
stepping pattern. Static metadata is supplied through dataclasses such as
\texttt{SurfaceSpec}, \texttt{ReactionWheelSpec},
\texttt{MagnetorquerSpec}, \texttt{ControlMomentGyroSpec}, and
\texttt{ThrusterSpec}. Runtime commands and sensor measurements live on the
data object, matching the model/data split familiar to MuJoCo users.

This paper will make four main contributions. First, it describes a
MuJoCo-compatible state and stepping interface for coupled orbital and
multibody simulation. Second, it gives the wrench-coupling formulation used
to apply orbital, environmental, and actuator effects to each body in the
LVLH frame. Third, it documents implementation details that are easy to get
wrong in MuJoCo spacecraft simulations, including body-frame angular velocity,
inertia-frame orientation, and world-frame external torques. Fourth, it
presents a validation and demonstration suite covering analytical orbital
limits, attitude dynamics, actuator behavior, articulated motion, and contact.

The remainder of the paper is organized as follows. Section~\ref{sec:related}
positions \mjorbit{} relative to astrodynamics frameworks and robotics
physics engines. Section~\ref{sec:formulation} introduces the coupled state,
frames, forces, and stepping algorithm. Section~\ref{sec:experiments}
outlines the validation studies and demonstrations planned for the submission.
Section~\ref{sec:conclusion} summarizes the intended contribution and open
development items.
